\section{Introduction}
\label{sec:intro}

Road traffic injuries constitute one of the leading causes of preventable death worldwide.
According to the World Health Organization, approximately \textbf{1.19 million people} die each year
as a result of road traffic crashes~\cite{WHO2023}.
In the United States alone, the National Highway Traffic Safety Administration reported that
distraction-affected crashes accounted for \textbf{3,275 fatalities} in 2022,
representing roughly 8\% of all traffic deaths that year~\cite{NHTSA2025}.
Despite decades of road safety improvements, distracted and cognitively impaired driving
remains a persistent and growing threat, exacerbated by the proliferation of in-vehicle
infotainment systems and mobile devices~\cite{Young2007, Kaplan2015}.

Existing approaches to driver impairment detection fall broadly into three categories,
each carrying fundamental practical limitations.
\textit{Physiological monitoring}---relying on heart rate, electrodermal activity, or
electroencephalography---offers direct access to the driver's arousal state but requires
intrusive wearable or embedded sensors that are costly, susceptible to motion artefacts,
and difficult to deploy at scale~\cite{Healey2005}.
\textit{Camera-based systems} that track eye gaze, head pose, or facial expression achieve
strong performance in controlled laboratory settings but raise privacy concerns and
degrade under variable lighting and partial occlusion~\cite{Vora2018}.
A third line of work frames the problem as \textit{crash or error prediction}---forecasting
discrete safety-critical events from sensor streams.
This formulation is, however, fundamentally ill-posed: errors are rare and highly
context-dependent, and their occurrence is jointly determined by the driver's state,
road geometry, and surrounding traffic---factors that in-vehicle sensors alone cannot
fully observe~\cite{Dingus2016}.
Naturalistic driving data confirm that driver-related factors are present in approximately
88\% of crashes, yet no deployed system reliably forecasts the \textit{next} error;
what remains feasible is detecting the driver \textit{state} that elevates error probability~\cite{Dingus2016}.

This distinction is not merely semantic.
A well-established body of work in driving psychology demonstrates that
vehicle kinematics carry reliable information about driver workload and attentional state.
Steering entropy theory, introduced by Nakayama et al.~\cite{Nakayama1999} and
further developed by Boer~\cite{Boer2000}, formalises this intuition: under increased
cognitive load, steering corrections become less predictable and more irregular,
reflecting degraded closed-loop vehicle control.
Subsequent work has confirmed that steering wheel reversal rates are sensitive to
cognitive secondary task load~\cite{Kountouriotis2016},
that cognitive distraction produces measurable degradation in speed maintenance and
braking response~\cite{Strayer2003, Strayer2006},
and that visual and cognitive distraction produce distinct, identifiable signatures
across speed regulation, lane-keeping precision, and lateral acceleration~\cite{Engstrom2005}.
A comprehensive theoretical account of these effects is provided by
Engstr\"{o}m et al.~\cite{Engstrom2017}, who establish that cognitive load
systematically reduces the bandwidth of the driver's closed-loop vehicle control,
producing observable kinematic consequences before overt errors occur.
Taken together, this body of evidence establishes that in-vehicle kinematic signals
carry systematic information about driver impairment state---a finding grounded in
decades of human factors research~\cite{Nakayama1999, Boer2000, Engstrom2017}.

Yet this established link has not been operationalised into what practitioners actually
need: a \textbf{fitness-to-drive index}---a continuous, real-time score that quantifies
how safely a driver is currently operating, without requiring physiological sensors,
cameras, or knowledge of the road environment.
Existing kinematic-based systems either rely on hand-crafted thresholds
that do not generalise across drivers, or are evaluated within-subject in ways that
overestimate real-world performance~\cite{Kaplan2015}.
We argue that the correct framing is not \textit{``will this driver make an error
in the next N seconds?''} --- a question that is ill-posed given the rarity and
context-dependence of errors~\cite{Dingus2016} --- but rather:
\textit{``is this driver currently in a state associated with elevated error risk?''}
The former demands event prediction; the latter demands \textbf{impairment state scoring},
which kinematics are well-suited to provide.

We present \textbf{KinTCN}, a Temporal Convolutional Network~\cite{Bai2018} trained on
four standard in-vehicle kinematic signals---speed, steering wheel angle, steering torque,
and lateral acceleration---collected during a controlled driving simulator study
($N = 30$~participants).
Safety-critical error events serve as \textit{weak supervision}: windows preceding
an error are used to define elevated-risk periods, but the model output is interpreted
not as an error probability but as a continuous \textbf{fitness-to-drive index}---a
well-calibrated score that ranks impaired driving periods above safe ones.
Evaluation is conducted under \textbf{Leave-One-Participant-Out Cross-Validation (LOPO-CV)},
the strictest available cross-participant protocol, ensuring that the index generalises
to drivers never seen during training.

The contributions of this work are threefold:

\begin{itemize}
    \item \textbf{A data-driven fitness-to-drive index.}
          We propose a continuous, well-calibrated impairment score derived exclusively
          from four standard in-vehicle kinematic signals, requiring no physiological
          sensors, cameras, or external infrastructure.
          The index is trained using safety-critical error events as weak supervision
          and is explicitly framed as a measure of driver fitness rather than
          a predictor of discrete crash events.

    \item \textbf{KinTCN: a temporal deep learning approach with rigorous evaluation.}
          We develop and evaluate KinTCN, a Temporal Convolutional Network that learns
          kinematic impairment signatures directly from raw time series.
          Evaluation is conducted under Leave-One-Participant-Out Cross-Validation,
          the strictest available cross-participant protocol, achieving
          AUROC~$= 0.781$ (ECE~$= 0.029$) and significantly outperforming competitive
          baselines (Wilcoxon $p = 0.0014$).
          Ablation studies over lookback duration and prediction horizon, together with
          a false positive rate audit on error-free participants, further validate the
          robustness of the approach.

    \item \textbf{A labelled driving simulator dataset.}
          We contribute a richly annotated dataset collected from $N = 30$ participants
          under controlled distraction conditions, comprising synchronised kinematic
          signals and timestamped safety-critical error events across multiple
          distraction types and route configurations.
\end{itemize}

The remainder of the paper is organised as follows.
Section~\ref{sec:related} reviews related work on driver impairment detection and
kinematic-based driver monitoring.
Section~\ref{sec:dataset} describes the simulator study and dataset.
Section~\ref{sec:method} details the KinTCN method and evaluation protocol.
Section~\ref{sec:results} presents experimental results.
Section~\ref{sec:discussion} discusses findings and limitations.
Section~\ref{sec:conclusion} concludes.
